\scanpage{37}
\mspage{35}

\deleted{\uline{Corollaire} Soit $f:X'\longrightarrow X^{\unsure{\prime|\;} }$ un foncteur qui soit une équivalence d'homotopie. \inserted{\unsure{Considérons} $M(X)\unsure{\relbar|\leftarrow}M(X')$, qui} Alors le foncteur correspondant $M(X^{\unsure{\prime|\;}})\Sigma_{X'}^{-1}\longrightarrow M$
\[
\mathrm{Hom}_{\Sigma_X}(a_X,b_X)\longrightarrow
\]}

Notons d'abord \deleted{que} (\unsure{ceci étant|c'était} implicite dans le Scholie) que toute flèche de Cat

\diagram{\begin{tikzcd}[ampersand replacement=\&]
 X \arrow[r,dash,"f"] \& Y
\end{tikzcd}}

définit

\diagram{\begin{tikzcd}[ampersand replacement=\&]
 M(Y)\arrow[r,"{M(f)}"] \& M(X)\quad\text{appliquant }\Sigma_Y\text{ dans }\Sigma_X,
\end{tikzcd}}

d'où

\[
 M(Y)\Sigma_Y^{-1}\overset{\mathrm{def}}{=}\mathbb D_{M,\Sigma}(Y)
 \longrightarrow M_X\Sigma_{\unsure{X|Y}}^{-1}=\mathbb D_{M,\Sigma}(X),
\]

\unsure{d'où}, pour tout $a,b\in\mathrm{Ob}\,M$, vu que $M(f)$ \deleted{applique} \inserted{transforme} $a_Y,\unsure{b_Y}$ en $a_X,b_X$ :

\[
 \boxed{\mathrm{Hom}_{\Sigma_Y}(a_Y,b_Y)\longrightarrow\mathrm{Hom}_{\Sigma_X}(a_X,b_X)}.
\]

Ceci posé, on a :

\uline{Corollaire} : Si $f$ est \unsure{une|un} \deleted{\unsure{équivalence d'}} \inserted{homotopisme}, alors quel que soient $a,b\in M$, la flèche précédente est un isomorphisme.

La question naturelle, c'est de savoir si la $=$ \deleted{\unsure{question}} conclusion est encore valable, en supposant seulement que $f$ soit une ``\uline{quasi-isomorphisme}'' (ou ``équivalence faible''), i.e que $f$ définisse un isom. dans la cat. \unsure{de|des} fractions Hot de Cat. (À ne pas confondre avec Cat hot,

% archive sheet number: 36
